\chapter{Weak-Axisymmetric Transport, Quotient Coordinates, and Orbit Closure}
\label{chap:part05_weak_axisymmetric_orbit}

\section*{Stage range: 164--200}
\addcontentsline{toc}{section}{Stage range: 164--200}

Part IV ended with the co-evolving core--mouth branch reduced to a first-order off-family normal coordinate.  Part V starts from that coordinate and asks a sharper question: after the Family--1 mouth/core compensation machinery is already in place, which microscopic grouped weak-axisymmetric motions still survive as genuine reduced defects, and which motions are only similarity-orbit relabelings of the same coherent branch?

The answer is that the apparent many-variable weak-axisymmetric problem collapses to two levels.  On the branch side, the grouped real \(P_2\) carrier reduces to a single transported prefactor slope,
\begin{equation}
\label{eq:part05-intro-xi1}
\Xi_1=\frac{P_1}{P_0},
\end{equation}
with the weak-axisymmetric grouped signature
\begin{equation}
\label{eq:part05-intro-lambda-signature}
\lambda_{20}=1,
\qquad
\lambda_{21}=\frac12,
\qquad
\lambda_{22}=-1.
\end{equation}
Equivalently, the grouped anisotropy lies on the fixed line
\begin{equation}
\label{eq:part05-intro-b3a}
b_{P_0}=3a_{P_0}.
\end{equation}
On the orbit side, the coherent local branch reduces to three exact quotient coordinates,
\begin{equation}
\label{eq:part05-intro-quotient-coordinates}
\mathfrak C_{{\rm tr},*},
\qquad
\mathfrak C_{{\rm nt},*},
\qquad
\epsilon_\eta,
\end{equation}
whose finite level sets are exactly the five-parameter microscopic similarity orbits.  The final reduced verdict after this part is therefore not an eight-slot branch residual plus an ambiguous orbit condition.  It is the exact four-scalar packet
\begin{equation}
\label{eq:part05-intro-final-packet}
\Delta_{\rm full}
=
(\Delta_Q,q_{\rm tr},q_{\rm nt},q_\eta),
\qquad
\Delta_Q:=\chi_Q-1.
\end{equation}

\section*{Part V at a glance}

\begin{readermap}
\item[Primary question] After Part IV reduces the retarded obstruction to one transported branch scalar, which weak-axisymmetric motions are genuine defects and which are only similarity-orbit relabelings of the same coherent branch?
\item[Starting inputs] The Part IV off-family normal coordinate \(\delta_\perp\), the compensated Family--1 mouth/core branch, and the outgoing-normalization packet.
\item[Delivers] The scalar transfer-shape slope \(\Xi_1=P_1/P_0\), the coherent normal form and monomial quotient coordinates, the Packet-A finish line \(\Delta_Q=\chi_Q-1\), and the final verdict packet \(\Delta_{\rm full}\).
\item[Used by] Part VI as the realization-compiler input and Parts VII--VIII as the strict and relaxed branch packet language.
\item[Result anchors] \resultanchor{MTDC-T9}.
\end{readermap}

\paragraph{Stage packets.}
\begin{center}
\small
\begin{tabular}{@{}>{\raggedright\arraybackslash}p{0.11\textwidth}>{\raggedright\arraybackslash}p{0.24\textwidth}>{\raggedright\arraybackslash}p{0.19\textwidth}>{\raggedright\arraybackslash}p{0.38\textwidth}@{}}
\toprule
Stages & Packet & Status & Carry-forward output \\
\midrule
164--169 & Off-family drift reduction and scalar-slippage firewall & \StatusExactClosure{} & Rewrites \(\delta_\perp\) into microscopic drifts, isolates the true scalar slippage channels, and proves pure grouped real \(P_2\) anisotropy does not linearly feed them. \\
\addlinespace
170--180 & Grouped outlet/load collapse to the transfer-shape scalar & \StatusExactClosure{} & Maps grouped outlet deformations to physical slopes, derives \(\Xi_{\rm load}\) and the static self-similarity laws, and collapses weak-axisymmetric outgoing transport to \(\Xi_1=P_1/P_0\). \\
\addlinespace
181--190 & Coherent normal form, monomial coordinates, and direct compilers & \StatusExactClosure{} & Compresses the coherent branch to \((\Sigma_{\rm tr},\Sigma_{\rm nt},\Sigma_\eta)\), identifies direct monomial quotient coordinates, and records observable compilers and no-go packets. \\
\addlinespace
191--197 & Packet-A reduction and retarded normalization finish line & \StatusExactClosure{} / \StatusReduced{} & Defines Packet A and Packet B, fixes the outgoing \(l=2\) fingerprint, imports the natural point-particle source map, and proves \(\Delta_Q=\chi_Q-1\) is the sole Packet-A scalar. \\
\addlinespace
198--200 & Orbit transport and final reference-free verdict packet & \StatusExactClosure{} & Removes representative privilege, proves exact orbit transport, and closes the home stretch with \(\Delta_{\rm full}=(\Delta_Q,q_{\rm tr},q_{\rm nt},q_\eta)\). \\
\bottomrule
\end{tabular}
\end{center}

\claimstatus{Mixed but predominantly \StatusExactClosure{}.  Stages 164--200 are exact algebra inside the declared co-evolving Family--1, grouped real \(P_2\), coherent local D/N, outgoing-prefactor, Packet-A, and Packet-B reduced hierarchies.  The natural point-particle source-map use in Stage 195 is a \StatusReduced{} import.  The actual nonlinear moving-throat PDE still has to return branch data that land on the final packet; realization by that branch remains \StatusOpen{} even though the compilers and quotient/orbit theorems are exact within closure.}

\derivationfield{What this part proves}{Part V reduces the apparent many-variable weak-axisymmetric defect problem to one transported branch scalar, three exact quotient coordinates, and the final four-scalar verdict packet \(\Delta_{\rm full}\).}
\derivationfield{What this part does not prove}{It does not show that the actual branch makes that packet vanish. It only proves that this is the exact reduced packet that later realization work must evaluate.}
\derivationfield{Why later parts need it}{Part VI turns \(\Delta_{\rm full}\) into the realization compiler, while Parts VII and VIII use the same packet language to state strict actual-branch tests and relaxed-branch extensions.}

\section{Block contract and theorem-level output}
\label{sec:part05-block-contract}

The contract of Part V is to turn the remaining first-order weak-axisymmetric and retarded branch checks into explicit finite packets.  It is not trying to prove that the completed nonlinear PDE realizes those packets.  Instead it proves that, once a candidate moving-throat solution gives the required reduced branch data, the verdict can be computed by exact formulas with no remaining interpretation freedom.

The block has two sides.  The first side is the weak-axisymmetric branch transport:
\begin{equation}
\label{eq:part05-chain-branch-side}
\delta_\perp
\longrightarrow
(\mathcal K_A,\mathcal G_A)
\longrightarrow
(u_2^{(1)},P_1/P_0)
\longrightarrow
\Xi_{\rm load}
\longrightarrow
\Xi_1.
\end{equation}
The second side is the coherent orbit quotient:
\begin{equation}
\label{eq:part05-chain-orbit-side}
(\Sigma_{\rm tr},\Sigma_{\rm nt},\Sigma_\eta)
\longrightarrow
(\mathfrak C_{{\rm tr},*},\mathfrak C_{{\rm nt},*},\epsilon_\eta)
\longrightarrow
\mathcal M_+/\mathcal G_*
\longrightarrow
(q_{\rm tr},q_{\rm nt},q_\eta).
\end{equation}
The final theorem joins these with the retarded Packet-A scalar \(\Delta_Q=\chi_Q-1\).

\begin{theorem}[Part V weak-axisymmetric quotient and orbit-closure theorem]
\label{thm:part05-main}
Inside the carried moving-throat reduced hierarchy for Stages 164--200, the following theorem chain holds.
\begin{enumerate}[leftmargin=2.2em]
  \item The first off-family Family--1 normal coordinate is an explicit logarithmic transport law in the moving-throat wall/BdG/Maxwell/mixed branch variables.  Arbitrary first-order isotropic bundle drift is tangent to the compensated parent family; true off-bundle motion first enters through three scalar slippages.
  \item A pure grouped real \(P_2\) anisotropy cannot linearly feed the scalar slippage channel.  Its linear effect enters the direct outlet coefficients only through
  \begin{equation}
  \label{eq:part05-theorem-KG}
  \mathcal K_A:=\delta D_{A,2}+\frac19\delta D_{A,0},
  \qquad
  \mathcal G_A:=\delta N_{A,0}-P_0\delta D_{A,0}.
  \end{equation}
  \item On the canonical compensated branch, these outlet coefficients are exactly the physical grouped slopes,
  \begin{equation}
  \label{eq:part05-theorem-physical-slopes}
  \mathfrak K_1=-D_0u_2^{(1)},
  \qquad
  \mathfrak G_1=D_0P_1.
  \end{equation}
  \item On the even-preserving weak-axisymmetric branch, all conservative grouped transport is controlled by one static operator slope, and the sole remaining linear grouped normalization defect is
  \begin{equation}
  \label{eq:part05-theorem-xiload}
  \Xi_{\rm load}
  :=
  \frac{N_{01}}{N_0}-\frac{D_{01}}{D_0}.
  \end{equation}
  \item The same scalar is the weighted failure of static self-similarity and, after wall-normalized outgoing-port factorization, the logarithmic slope of the effective transfer shape:
  \begin{equation}
  \label{eq:part05-theorem-xi1-transfer}
  \Xi_1
  =
  \frac{P_1}{P_0}
  =
  \frac{\delta\ln\mathcal T_{{\rm eff},A}^{2}}{\epsilon\lambda_A}.
  \end{equation}
  \item On the coherent local D/N branch, the weak-axisymmetric observable drifts obey the triangular normal form
  \begin{equation}
  \label{eq:part05-theorem-triangular}
  \Theta_1=-C_{\rm tr}\Sigma_{\rm tr},
  \qquad
  \Xi_1=A_{\rm tr}\Sigma_{\rm tr}+\Sigma_{\rm nt},
  \qquad
  \mathcal R_1+\Xi_1=-\frac{\epsilon_\eta}{1-\epsilon_\eta}\Sigma_\eta.
  \end{equation}
  \item The three normal-form coordinates are exactly the logarithmic drifts of three direct microscopic monomials,
  \begin{equation}
  \label{eq:part05-theorem-monomials}
  \delta\ln\mathfrak C_{{\rm tr},*}=\Sigma_{\rm tr},
  \qquad
  \delta\ln\mathfrak C_{{\rm nt},*}=\Sigma_{\rm nt},
  \qquad
  \delta\ln\epsilon_\eta=\Sigma_\eta.
  \end{equation}
  \item The finite level sets of these three monomials are exactly the five-parameter similarity orbits \(\mathcal G_*\), so
  \begin{equation}
  \label{eq:part05-theorem-quotient}
  \mathcal M_+/\mathcal G_*
  \cong
  (\mathbb R_{>0})^3,
  \qquad
  (\mathfrak C_{{\rm tr},*},\mathfrak C_{{\rm nt},*},\epsilon_\eta)
  \text{ are quotient coordinates.}
  \end{equation}
  \item The branch-side Packet A collapses to one scalar finish line,
  \begin{equation}
  \label{eq:part05-theorem-packetA}
  \Delta_{\rm branch}=0
  \iff
  \chi_Q=1
  \iff
  \Delta_Q:=\chi_Q-1=0.
  \end{equation}
  \item The full reference-free reduced verdict is therefore the four-scalar packet
  \begin{equation}
  \label{eq:part05-theorem-finalpacket}
  \Delta_{\rm full}^{(x\mid\mathcal O_*)}
  :=
  \bigl(
  \Delta_Q(x),
  q_{\rm tr}^{(x\leftarrow\mathcal O_*)},
  q_{\rm nt}^{(x\leftarrow\mathcal O_*)},
  q_\eta^{(x\leftarrow\mathcal O_*)}
  \bigr),
  \end{equation}
  with
  \begin{equation}
  \label{eq:part05-theorem-finalzero}
  \Delta_{\rm full}^{(x\mid\mathcal O_*)}=0
  \iff
  \Delta_{\rm branch}(x)=0
  \text{ and }
  x\in\mathcal O_*.
  \end{equation}
\end{enumerate}
\end{theorem}

\begin{proof}[Proof structure]
Stages 164--168 translate the first off-family normal coordinate into explicit microscopic transport variables and then separate on-bundle tangent motion from off-bundle slippage.  Stages 169--173 prove that grouped real \(P_2\) anisotropy does not linearly feed scalar slippage, derive the direct outlet map, and collapse the even-preserving grouped defect to \(\Xi_{\rm load}\).  Stages 174--180 factor \(\Xi_{\rm load}\) through wall-referenced self-similarity, outgoing-load factors, port co-loading, and the effective transfer shape, giving \(\Xi_1=P_1/P_0\).  Stages 181--185 pass from physical coherent-kernel variables to the triangular normal form and then to direct microscopic monomials.  Stages 186--187 prove the finite orbit--fibre theorem.  Stages 188--192 build the observable, transfer-shape, finite-packet, and projector compilers.  Stages 193--197 close the branch-side Packet-A finish line to \(\Delta_Q=\chi_Q-1\).  Stages 198--200 remove the orbit representative and combine Packet A and Packet B into \eqref{eq:part05-theorem-finalpacket}.
\end{proof}

\section{Weak-axisymmetric grouped signature}
\label{sec:part05_weak_axisymmetric_orbit:weak-axisymmetric-grouped-signature}

\subsection{Local goal}
The first job in Part V is to connect the co-evolving Family--1 off-family scalar from Part IV to the grouped real \(P_2\) transport language.  Stage 163 had left the first true normal coordinate in the form
\begin{equation}
\label{eq:part05-delta-perp-original}
\delta_\perp
=
\mathfrak g_*
\delta\ln\!\left(\frac{g_qK_s}{g_s\lambda}\right)
+
\frac{1}{4\sqrt{1+\mathfrak r_*^2}}
\delta\ln\!\left(\frac{K_sK_q}{\lambda^2}\right).
\end{equation}
Stage 164 shows that this is not an abstract parent-coordinate expression.  With
\begin{equation}
\label{eq:part05-rg-defs}
\mathfrak r=\frac{\lambda}{\sqrt{K_sK_q}},
\qquad
\mathfrak g=\frac{g_q\sqrt{K_s}}{g_s\sqrt{K_q}},
\qquad
r_c=\frac{\lambda^2}{K_sK_q}=\mathfrak r^2,
\end{equation}
the two channels are exactly
\begin{equation}
\label{eq:part05-log-channels-rg}
\delta\ln\!\left(\frac{g_qK_s}{g_s\lambda}\right)
=
\delta\ln\!\left(\frac{\mathfrak g}{\mathfrak r}\right),
\qquad
\delta\ln\!\left(\frac{K_sK_q}{\lambda^2}\right)
=-\delta\ln r_c.
\end{equation}
Thus \(\delta_\perp\) measures relative mouth-coupling drift and hybridization drift.

\subsection{Explicit branch-variable form}
On the healing-locked wall/BdG/Maxwell/mixed branch, the same scalar becomes
\begin{equation}
\label{eq:part05-delta-perp-explicit}
\begin{aligned}
\delta_\perp={}&
A_*\,\delta\ln\!\left(\frac{\mathcal Z_q}{\rho_w}\right)
+3A_*\,\delta\ln c_{s,w}
+B_*\,\delta\ln c_s
-\mathfrak g_*\,\delta\ln\mathcal T_m \\
&-(\mathfrak g_*+B_*)\delta\ln v_{w0}
-2A_*\,\delta\ln a
-C_*\,\delta\ln L_W,
\end{aligned}
\end{equation}
where
\begin{equation}
\label{eq:part05-ABC-defs}
A_*:=\mathfrak g_*+\frac{1}{4\sqrt{1+\mathfrak r_*^2}},
\qquad
B_*:=\frac{1}{2\sqrt{1+\mathfrak r_*^2}},
\qquad
C_*:=2\mathfrak g_*+\frac{3}{4\sqrt{1+\mathfrak r_*^2}}.
\end{equation}
The exact tangency condition \(\delta_\perp=0\) is therefore a co-transport law for \(\mathcal T_m\) relative to
\(
(\mathcal Z_q,\allowbreak \rho_w,\allowbreak c_{s,w},\allowbreak c_s,\allowbreak v_{w0},\allowbreak a,\allowbreak L_W)
\),
not a free adjustment.

Stages 165--167 then show that the explicit lower compensated branch is even more constrained.  The D/N tube length co-transports with the mouth radius,
\begin{equation}
\label{eq:part05-LW-co-transport}
\delta\ln L_W=\delta\ln a,
\end{equation}
and the exact branch laws reduce the apparent branch drift problem to four irreducible variables.  Stage 166 chooses the four bundle observables
\begin{equation}
\label{eq:part05-bundle-observables}
\Theta_w,
\qquad
K_s,
\qquad
K_q,
\qquad
P_0=\frac{N_0}{D_0}
\end{equation}
as the right variables to invert.  Stage 167 proves that arbitrary first-order drift of these four bundle observables leaves the compensation ratios tangent to the parent family:
\begin{equation}
\label{eq:part05-tangent-parent-family}
\delta\ln r_c=0,
\qquad
\delta\ln\mathfrak r=0,
\qquad
\delta\ln\mathfrak g=0
\quad
\text{for pure bundle drift.}
\end{equation}
Consequently, any first-order failure after Stage 167 must be genuinely off-bundle.

\subsection{Off-bundle slippages and grouped signature}
Stage 168 introduces the three scalar slippages
\begin{equation}
\label{eq:part05-slippage-names}
\varepsilon_L,
\qquad
\varepsilon_v,
\qquad
\varepsilon_T,
\end{equation}
which measure failure of the exact lower-branch laws for axial length, mixed background speed, and mouth traction.  Stage 169 then proves a key firewall: a pure weak grouped real \(P_2\) anisotropy cannot linearly source these scalar slippages.  The reason is angular.  The weak grouped perturbation has zero weighted trace at first order; scalar feed-down can only be built from quadratic grouped invariants such as
\begin{equation}
\label{eq:part05-grouped-invariant}
\mathcal I[x,y]
=
\frac15\delta x^T G_{\rm grp}\delta y
=
4a_xa_y+\frac45 b_xb_y,
\qquad
G_{\rm grp}=\operatorname{diag}(1,2,2).
\end{equation}
Thus linear grouped real \(P_2\) transport enters the direct outlet coefficients rather than the scalar off-bundle slippage channel.

The grouped weak-axisymmetric signature is inherited from the first axisymmetric quadrupole perturbation:
\begin{equation}
\label{eq:part05-weak-axisymmetric-signature}
\lambda_{20}=1,
\qquad
\lambda_{21}=\frac12,
\qquad
\lambda_{22}=-1.
\end{equation}
For any first-order grouped quantity
\begin{equation}
\label{eq:part05-lambda-expansion-general}
x_A=x_0+\epsilon\lambda_Ax_1+O(\epsilon^2),
\end{equation}
the weighted grouped anomalies obey
\begin{equation}
\label{eq:part05-b-equals-3a-general}
a_x=\frac{\epsilon}{4}x_1,
\qquad
b_x=\frac{3\epsilon}{4}x_1,
\qquad
b_x=3a_x.
\end{equation}
This fixed line is the weak-axisymmetric branch signature used throughout the rest of Part V.

\claimstatus{Stages 164--169 are \StatusExactClosure{} inside the explicit co-evolving Family--1 throat-core closure and the grouped real \(P_2\) weak-axisymmetric expansion.  They do not prove that the actual nonlinear PDE produces a particular slippage; they identify the exact variables through which such slippage would appear.}

\section{Even-preserving and hidden-even gates}
\label{sec:part05_weak_axisymmetric_orbit:even-preserving-and-hidden-even-gates}

\subsection{Direct outlet map}
Stages 170--171 move the linear grouped problem from scalar slippage to direct outlet data.  For each grouped lane \(A\), define
\begin{equation}
\label{eq:part05-KG-defs-section}
\mathcal K_A:=\delta D_{A,2}+\frac19\delta D_{A,0},
\qquad
\mathcal G_A:=\delta N_{A,0}-P_0\delta D_{A,0}.
\end{equation}
The hidden even deformation and hidden odd normalization deformation are then
\begin{equation}
\label{eq:part05-dk-dg-outlet}
\delta\kappa_W^{(A)}
=
\frac{3(1-\sigma_*)}{\sigma_*D_0}\mathcal K_A,
\qquad
\delta\gamma_W^{(A)}
=-\frac{1-\sigma_*}{9\sigma_*N_0}\mathcal G_A.
\end{equation}
Stage 171 decomposes these into microscopic sectors:
\begin{equation}
\label{eq:part05-K-decomp}
\mathcal K_A
=
\mathcal W_A-\mathcal B_A-\mathcal Z_A,
\end{equation}
with
\begin{equation}
\label{eq:part05-WBZ-defs}
\mathcal W_A:=\frac19\delta K_A-\delta M_A,
\qquad
\mathcal B_A:=\delta B_{A,2}+\frac19\delta B_{A,0},
\qquad
\mathcal Z_A:=\delta Z_{A,2}+\frac19\delta Z_{A,0},
\end{equation}
and
\begin{equation}
\label{eq:part05-G-decomp}
\mathcal G_A
=
-P_0\delta K_A+P_0\delta B_{A,0}+\mathcal N_A,
\qquad
\mathcal N_A:=\delta N_{A,0}+P_0\delta Z_{A,0}.
\end{equation}
This identifies the only four microscopic bundles that can generate a linear grouped outlet obstruction: wall baseline, BdG support, conservative Maxwell/mixed, and outgoing-transfer loading.

\subsection{Physical slope collapse}
Stage 172 rewrites the same obstruction in physical grouped-response variables.  For each grouped lane,
\begin{equation}
\label{eq:part05-u2-u4-P0-defs}
u_2^{(A)}=-\frac{D_{A,2}}{D_{A,0}},
\qquad
u_4^{(A)}=\frac{D_{A,2}^2-D_{A,0}D_{A,4}}{D_{A,0}^2},
\qquad
P_0^{(A)}=\frac{N_{A,0}}{D_{A,0}}.
\end{equation}
On the canonical compensated branch, the two microscopic amplitudes are simply
\begin{equation}
\label{eq:part05-KG-physical}
\mathfrak K_1=-D_0u_2^{(1)},
\qquad
\mathfrak G_1=D_0P_1.
\end{equation}
Equivalently,
\begin{equation}
\label{eq:part05-dk-dg-physical}
\delta\kappa_W^{(A)}
=-\frac{3(1-\sigma_*)}{\sigma_*}\delta u_2^{(A)},
\qquad
\delta\gamma_W^{(A)}
=-\frac{1-\sigma_*}{9\sigma_*}\frac{\delta P_0^{(A)}}{P_0}.
\end{equation}
So the full linear grouped outlet problem is now measurable by \(u_2^{(1)}\) and \(P_1/P_0\).

\subsection{Even-preserving branch}
Stage 173 then computes those physical slopes directly from the weak-axisymmetric grouped operator expansion
\begin{equation}
\label{eq:part05-DN-expansion}
D_{A,0}=D_0+\epsilon\lambda_AD_{01}+O(\epsilon^2),
\quad
D_{A,2}=D_2+\epsilon\lambda_AD_{21}+O(\epsilon^2),
\quad
D_{A,4}=D_4+\epsilon\lambda_AD_{41}+O(\epsilon^2),
\end{equation}
\begin{equation}
\label{eq:part05-N-expansion}
N_{A,0}=N_0+\epsilon\lambda_AN_{01}+O(\epsilon^2).
\end{equation}
The exact slope formulas are
\begin{equation}
\label{eq:part05-physical-slopes-stage173}
u_2^{(1)}=-\frac{D_{21}+u_2D_{01}}{D_0},
\end{equation}
\begin{equation}
\label{eq:part05-u4-slope-stage173}
u_4^{(1)}
=-\frac{5D_{01}+18D_{21}+81D_{41}}{81D_0}
\quad
\text{on }(u_2,u_4)=\left(\frac19,\frac{4}{81}\right),
\end{equation}
and
\begin{equation}
\label{eq:part05-P1P0-stage173}
\frac{P_1}{P_0}
=
\frac{N_{01}}{N_0}-\frac{D_{01}}{D_0}.
\end{equation}
If the canonical conservative even fingerprint is preserved at first order, then
\begin{equation}
\label{eq:part05-even-preserving-D21}
D_{21}=-\frac{D_{01}}{9}.
\end{equation}
The hidden-even consistency relation \(u_4^{(1)}=(8/9)u_2^{(1)}\) is exactly
\begin{equation}
\label{eq:part05-hidden-even-D41-general}
D_{41}=\frac23D_{21}+\frac{1}{27}D_{01},
\end{equation}
and therefore on the even-preserving branch
\begin{equation}
\label{eq:part05-even-preserving-D41}
D_{41}=-\frac{D_{01}}{27}.
\end{equation}
Thus the conservative grouped response is frozen to canonical order, and the only surviving linear grouped defect is the static loading mismatch
\begin{equation}
\label{eq:part05-Xiload-def}
\Xi_{\rm load}
:=
\frac{N_{01}}{N_0}-\frac{D_{01}}{D_0}.
\end{equation}
Its lane form is
\begin{equation}
\label{eq:part05-Xiload-lanes}
\Delta_Q^{(20)}=\epsilon\Xi_{\rm load},
\qquad
\Delta_Q^{(21)}=\frac{\epsilon}{2}\Xi_{\rm load},
\qquad
\Delta_Q^{(22)}=-\epsilon\Xi_{\rm load}.
\end{equation}

\claimstatus{Stages 170--173 are \StatusExactClosure{} in the grouped low-frequency bundle.  The even-preserving condition is a branch restriction, not an unconditional statement that the completed PDE must satisfy it.}

\section{Outgoing-load scalar \texorpdfstring{\((\Xi_1=P_1/P_0)\)}{(Xi1=P1/P0)}}
\label{sec:part05_weak_axisymmetric_orbit:outgoing-load-scalar-xi-1-p-1-p-0}

\subsection{Self-similarity form of the loading mismatch}
Stage 174 rewrites the mismatch \eqref{eq:part05-Xiload-def} as failure of static self-similarity.  Let
\begin{equation}
\label{eq:part05-static-slope-defs}
\delta_K:=\frac{K_1}{K},
\qquad
\delta_B:=\frac{B_0^{(1)}}{B_0},
\qquad
\delta_Z:=\frac{Z_0^{(1)}}{Z_0},
\qquad
\delta_N:=\frac{N_1}{N_0},
\end{equation}
and
\begin{equation}
\label{eq:part05-weights-static}
\omega_B:=\frac{B_0}{D_0},
\qquad
\omega_Z:=\frac{Z_0}{D_0},
\qquad
D_0=K-B_0-Z_0.
\end{equation}
Then
\begin{equation}
\label{eq:part05-Xiload-selfsimilarity}
\Xi_{\rm load}
=
(\delta_N-\delta_K)
+\omega_B(\delta_B-\delta_K)
+\omega_Z(\delta_Z-\delta_K).
\end{equation}
Define the wall-referenced self-similarity defect fields
\begin{equation}
\label{eq:part05-Sigma-self-fields}
\Sigma_\alpha^{(B)}:=2\delta\ln\!\left(\frac{c_\alpha}{\varpi_\alpha}\right)-\delta_K,
\quad
\Sigma_r^{(Z)}:=\delta\ln\!\left(\frac{Q_r}{\Delta_r}\right)-\delta_K,
\quad
\Sigma_r^{(N)}:=2\delta\ln\!\left(\frac{P_r}{\Delta_r}\right)-\delta_K.
\end{equation}
Then
\begin{equation}
\label{eq:part05-Xiload-Sigma-fields}
\Xi_{\rm load}
=
\sum_r\rho_r^{(N)}\Sigma_r^{(N)}
+
\omega_B\sum_\alpha\rho_\alpha^{(B)}\Sigma_\alpha^{(B)}
+
\omega_Z\sum_r\rho_r^{(Z)}\Sigma_r^{(Z)}.
\end{equation}
Static self-similarity is the sufficient condition that all three defect families vanish.  On the even-preserving branch it implies \(\Xi_{\rm load}=0\) and hence no first-order grouped normalization defect.

\subsection{Outgoing-load theorem and mixed-leg law}
Stages 175--176 factor the port objects relative to the wall baseline.  The remaining defect is controlled by the outgoing load factor
\begin{equation}
\label{eq:part05-load-factor-Lambda}
\Lambda_r:=\frac{P_r}{\Delta_r}.
\end{equation}
On conservative-shape-preserving branches the full remaining defect is carried only by the outgoing load slope relative to the wall.  Stage 176 factors that load into a mixed leg, an interference ratio, and a hybridization ratio.  In the notation of the reduced one-port block,
\begin{equation}
\label{eq:part05-load-factor-factorization}
\frac{\Lambda_r^2}{K}
=
\mathcal M_r^2
\frac{(1+\mathcal I_r)^2}{(1-\mathcal H_r)^2}.
\end{equation}
Consequently
\begin{equation}
\label{eq:part05-SigmaN-factorized}
\Sigma_r^{(N)}
=
2\delta\ln\mathcal M_r
+
\frac{2\mathcal I_r}{1+\mathcal I_r}\delta\ln\mathcal I_r
+
\frac{2\mathcal H_r}{1-\mathcal H_r}\delta\ln\mathcal H_r.
\end{equation}
If the interference and hybridization ratios are rigid, the whole outgoing defect is the raw mixed-leg drift.  The zero-defect condition becomes the square-root mixed-leg law
\begin{equation}
\label{eq:part05-square-root-mixed-leg}
\frac{G_W}{\Omega_W^2}\propto\sqrt K.
\end{equation}
This is useful because it identifies a concrete microscopic path to zero defect rather than just a cancellation among compiled coefficients.

\subsection{Collapse to one scalar amplitude}
Stage 177 applies the weak-axisymmetric signature \eqref{eq:part05-weak-axisymmetric-signature}.  Every microscopic outgoing slippage inherits the same lane pattern, so the full remaining grouped defect is one scalar amplitude:
\begin{equation}
\label{eq:part05-Xi1-definition}
\Xi_{\rm load}^{(A)}=\epsilon\lambda_A\Xi_1.
\end{equation}
The same scalar is the physical outgoing-prefactor slope:
\begin{equation}
\label{eq:part05-Xi1-P1P0}
\Xi_1=\frac{P_1}{P_0}.
\end{equation}
Stage 178 rewrites it portwise.  For
\begin{equation}
\label{eq:part05-port-N-static}
N_{A,0}^{(r)}=\frac{P_{A,r}^2}{\Delta_{A,r}^2},
\end{equation}
define the port slope \(\nu_r\) by
\begin{equation}
\label{eq:part05-nu-r-def}
\delta\ln N_{A,0}^{(r)}=\epsilon\lambda_A\nu_r.
\end{equation}
Then
\begin{equation}
\label{eq:part05-Xi1-port-coloading}
\Xi_1=\bar\nu_N-\kappa_1,
\end{equation}
where \(\kappa_1\) is the wall-baseline static slope and \(\bar\nu_N\) is the outgoing-weighted average.  Thus zero defect is exact co-loading of outgoing transfer with the wall baseline.

Stages 179--180 put the same statement into transfer-shape language.  Each port factors as
\begin{equation}
\label{eq:part05-port-transfer-shape}
N_{A,0}^{(r)}=K_A\mathcal T_{A,r}^2,
\qquad
\mathcal T_{A,r}
=
\frac{\widehat G_{W,A,r}+\widehat R_{A,r}\widehat G_{U,A,r}}{1-\widehat R_{A,r}^2}.
\end{equation}
Therefore
\begin{equation}
\label{eq:part05-Xi1-transfer-shape-sum}
\Xi_1=2\sum_r\rho_r^{(N)}\tau_r,
\qquad
\tau_r:=\frac{\delta\ln\mathcal T_{A,r}}{\epsilon\lambda_A}.
\end{equation}
Equivalently, all ports collapse to one effective transfer shape,
\begin{equation}
\label{eq:part05-Teff-def}
\mathcal T_{{\rm eff},A}^2
:=
\sum_r\mathcal T_{A,r}^2
=
\frac{N_{A,0}}{K_A},
\end{equation}
and
\begin{equation}
\label{eq:part05-Xi1-Teff}
\Xi_1
=
\frac{\delta\ln\mathcal T_{{\rm eff},A}^2}{\epsilon\lambda_A}.
\end{equation}
On the actual minimal one-port continuum branch,
\begin{equation}
\label{eq:part05-T-one-port-continuum}
\mathcal T_A^2
=
\frac{Z_{W,A}(1+\rho_A)^2}{\Omega_{W,A}^2(1-\epsilon_{W,A})^2}
=
\frac{27\pi^2Gc_s^5}{20a^5c^5}
\frac{1-\epsilon_{\eta,A}}{R_{{\rm target},A}}.
\end{equation}
So the remaining theorem gate is to determine whether the actual weak-axisymmetric branch keeps \(\mathcal T_{{\rm eff}}\) rigid.

\claimstatus{Stages 174--180 are \StatusExactClosure{} inside the weak-axisymmetric grouped bundle and the declared outgoing-port factorization.  The square-root mixed-leg law is a sufficient zero-defect branch condition under interference/hybridization rigidity; actual realization is not assumed.}

\section{Coherent-kernel slippages}
\label{sec:part05_weak_axisymmetric_orbit:coherent-kernel-slippages}

\subsection{Actual coherent local D/N branch}
Stage 181 inserts the coherent local D/N tracking branch into \eqref{eq:part05-T-one-port-continuum}.  On that branch,
\begin{equation}
\label{eq:part05-coherent-T}
\mathcal T^2
=
\frac{Z_W(1+\chi_0)^2}{\Omega_W^2(1-\epsilon)^2}
=
\frac{27\pi^2Gc_s^5}{20a^5c^5}\frac{1-\epsilon_\eta}{R_{\rm target}},
\end{equation}
where
\begin{equation}
\label{eq:part05-epsilon-split}
\epsilon
=
\epsilon_W\left(1-\frac{2}{11}\frac{\delta_U}{1+\delta_U}\right).
\end{equation}
The weak-axisymmetric defect is
\begin{equation}
\label{eq:part05-Xi1-coherent-physical}
\Xi_1
=
\zeta_Z-\omega_W
+
\frac{2\chi_1}{1+\chi_0}
+
\frac{2\epsilon_1}{1-\epsilon},
\end{equation}
with
\begin{equation}
\label{eq:part05-epsilon1-coherent}
\epsilon_1
=
\left(1-\frac{2}{11}\frac{\delta_U}{1+\delta_U}\right)\varepsilon_W
-
\frac{2\epsilon_W}{11(1+\delta_U)^2}\delta_{U,1}.
\end{equation}
The support enhancement ratio \(\zeta\) drops out identically:
\begin{equation}
\label{eq:part05-support-blindness}
\partial_\zeta\ln\mathcal T^2=0,
\qquad
\partial_\zeta\ln R_{\rm target}=0,
\qquad
\partial_\zeta\Xi_1=0.
\end{equation}
Thus coherent support can help the steady normalization load, but it cannot repair the weak-axisymmetric prefactor slope at first order.

\subsection{Microscopic slippage decomposition}
Stage 182 pushes \eqref{eq:part05-Xi1-coherent-physical} down to microscopic coherent-kernel slippages.  The direct grouped defect depends on four slippages,
\begin{equation}
\label{eq:part05-slippage-four}
\Sigma_Z,
\qquad
\Sigma_\chi,
\qquad
\Sigma_\epsilon,
\qquad
\Sigma_\delta,
\end{equation}
and selected-branch dressing introduces one more,
\begin{equation}
\label{eq:part05-slippage-eta}
\Sigma_\eta.
\end{equation}
The direct defect law is
\begin{equation}
\label{eq:part05-Xi1-slippage-law}
\Xi_1
=
\Sigma_Z
+
\frac{2\chi_0}{1+\chi_0}\Sigma_\chi
+
\frac{2\epsilon_W}{1-\epsilon}
\left[
\frac{11+9\delta_U}{11(1+\delta_U)}\Sigma_\epsilon
-
\frac{2\delta_U}{11(1+\delta_U)^2}\Sigma_\delta
\right].
\end{equation}
The tracking drift is the special combination
\begin{equation}
\label{eq:part05-Sigma-tr-def}
\Sigma_{\rm tr}:=(1+\chi_0)\Sigma_\delta+(1+\delta_U)\Sigma_\chi,
\end{equation}
which controls the tracking observable
\begin{equation}
\label{eq:part05-Theta-Sigma-tr}
\Theta_1
=
-
\frac{\chi_0\delta_U}{(1+\chi_0)(1+\delta_U)(1+\chi_0+\delta_U)}
\Sigma_{\rm tr}.
\end{equation}
So exact tracking rigidity means \(\Sigma_{\rm tr}=0\), but that alone does not generally force \(\Xi_1=0\); a genuine nontracking transfer-shape slippage can remain.

\claimstatus{Stages 181--182 are \StatusExactClosure{} for the coherent local D/N tracking branch.  The support-blindness statement is exact inside that branch; it does not assert that support is irrelevant to all other branch tests.}

\section{Triangular normal form}
\label{sec:part05_weak_axisymmetric_orbit:triangular-normal-form}

\subsection{Normal-form coordinates}
Stage 183 packages the five slippages into three branch-adapted coordinates.  The nontracking transfer-shape slippage is defined by
\begin{equation}
\label{eq:part05-Sigma-nt-def}
\begin{aligned}
\Sigma_{\rm nt}:={}&
\Sigma_Z
+
\frac{2\epsilon_W}{1-\epsilon}\frac{11+9\delta_U}{11(1+\delta_U)}\Sigma_\epsilon \\
&-
\left[
\frac{2\chi_0}{1+\delta_U}
+
\frac{4\epsilon_W\delta_U}{11(1-\epsilon)(1+\delta_U)^2}
\right]\Sigma_\delta.
\end{aligned}
\end{equation}
The three branch-adapted coordinates are therefore
\begin{equation}
\label{eq:part05-normal-form-coordinates}
(\Sigma_{\rm tr},\Sigma_{\rm nt},\Sigma_\eta).
\end{equation}
In these coordinates the observable drifts have the exact triangular form
\begin{equation}
\label{eq:part05-triangular-normal-form}
\Theta_1=-C_{\rm tr}\Sigma_{\rm tr},
\qquad
\Xi_1=A_{\rm tr}\Sigma_{\rm tr}+\Sigma_{\rm nt},
\qquad
\mathcal R_1+\Xi_1=-\frac{\epsilon_\eta}{1-\epsilon_\eta}\Sigma_\eta,
\end{equation}
where
\begin{equation}
\label{eq:part05-Ctr-Atr-defs}
C_{\rm tr}
=
\frac{\chi_0\delta_U}{(1+\chi_0)(1+\delta_U)(1+\chi_0+\delta_U)},
\qquad
A_{\rm tr}
=
\frac{2\chi_0}{(1+\chi_0)(1+\delta_U)}.
\end{equation}
This is the first exact three-scalar normal form of the weak-axisymmetric coherent branch.

\subsection{Inversion and interpretation}
The triangular structure makes the three physical failure modes transparent.
\begin{enumerate}[leftmargin=2.2em]
  \item \(\Sigma_{\rm tr}\) is tracking failure.  It is measured directly by \(\Theta_1\).
  \item \(\Sigma_{\rm nt}\) is nontracking transfer-shape failure.  It is the residual part of \(\Xi_1\) after subtracting tracking feed-through.
  \item \(\Sigma_\eta\) is dressing failure.  It measures the mismatch between direct transfer-shape drift and selected-branch demand through \(\mathcal R_1+\Xi_1\).
\end{enumerate}
Zero defect in this normal form is
\begin{equation}
\label{eq:part05-zero-normal-form}
\Theta_1=\Xi_1=\mathcal R_1=0
\iff
\Sigma_{\rm tr}=\Sigma_{\rm nt}=\Sigma_\eta=0.
\end{equation}

\subsection{Branch-observable coordinates}
Stage 184 identifies exact branch composites whose logarithmic drifts are the same three coordinates.  The tracking factor is
\begin{equation}
\label{eq:part05-Rtr-def}
R_{\rm tr}
=
\frac{1+\chi_0/(1+\delta_U)}{1+\chi_0}
=
\frac{1+\chi_0+\delta_U}{(1+\chi_0)(1+\delta_U)}.
\end{equation}
The corrected nontracking composite is
\begin{equation}
\label{eq:part05-Nstar-def}
\mathfrak N_*:=\mathcal T^2R_{\rm tr}^{B_*},
\qquad
B_*:=\frac{2(1+\chi_{0,*}+\delta_{U,*})}{\delta_{U,*}}.
\end{equation}
After a harmless positive normalization of the tracking factor,
\begin{equation}
\label{eq:part05-Tstar-def}
\mathfrak T_*:=R_{\rm tr}^{-C_*},
\qquad
C_*:=\frac{(1+\chi_{0,*})(1+\delta_{U,*})(1+\chi_{0,*}+\delta_{U,*})}{\chi_{0,*}\delta_{U,*}},
\end{equation}
the branch-coordinate laws are
\begin{equation}
\label{eq:part05-branch-coordinate-laws}
\delta\ln\mathfrak T_* = \Sigma_{\rm tr},
\qquad
\delta\ln\mathfrak N_* = \Sigma_{\rm nt},
\qquad
\delta\ln\epsilon_\eta=\Sigma_\eta.
\end{equation}
So the next continuation point is no longer a list of microscopic slippages; it is drift of three exact branch composites.

\claimstatus{Stages 183--184 are \StatusExactClosure{} inside the coherent local D/N branch.  The normal form is a compiler from reduced coherent variables to observables; it does not assert that the completed PDE has zero coordinates.}

\section{Microscopic monomial invariants}
\label{sec:part05_weak_axisymmetric_orbit:microscopic-monomial-invariants}

\subsection{Direct monomial coordinates}
Stage 185 pushes the three branch composites to direct microscopic monomials.  At reference-branch order,
\begin{equation}
\label{eq:part05-monomial-drift-laws}
\delta\ln\mathfrak C_{{\rm tr},*}=\Sigma_{\rm tr},
\qquad
\delta\ln\mathfrak C_{{\rm nt},*}=\Sigma_{\rm nt},
\qquad
\delta\ln\epsilon_\eta=\Sigma_\eta.
\end{equation}
The tracking monomial is
\begin{equation}
\label{eq:part05-Ctr-monomial-simple}
\mathfrak C_{{\rm tr},*}
:=
\chi_0^{1+\delta_{U,*}}\delta_U^{1+\chi_{0,*}},
\end{equation}
and the nontracking monomial is
\begin{equation}
\label{eq:part05-Cnt-monomial-simple}
\mathfrak C_{{\rm nt},*}
:=
\frac{Z_W}{\Omega_W^2}\epsilon_W^{E_*}\delta_U^{-F_*},
\end{equation}
with
\begin{equation}
\label{eq:part05-EF-defs}
E_*=
\frac{2\epsilon_{W,*}}{1-\epsilon_*}\frac{11+9\delta_{U,*}}{11(1+\delta_{U,*})},
\qquad
F_*=
\frac{2\chi_{0,*}}{1+\delta_{U,*}}
+
\frac{4\epsilon_{W,*}\delta_{U,*}}{11(1-\epsilon_*)(1+\delta_{U,*})^2}.
\end{equation}
Using the coherent definitions, these become direct microscopic monomials in the positive variables
\begin{equation}
\label{eq:part05-micro-state-vector}
\mathbf x
=
(\lambda_W,
 c_{\eta U},
 \gamma,
 K_U,
 K_\eta^{({\rm eff})},
 K_W^{({\rm eff})},
 \mu_W,
 T_U).
\end{equation}
Explicitly,
\begin{equation}
\label{eq:part05-Ctr-micro}
\mathfrak C_{{\rm tr},*}
=
\left(\frac{\gamma c_{\eta U}}{K_U}\right)^{1+\delta_{U,*}}
\left(\frac{\pi^2T_U}{L^2K_U}\right)^{1+\chi_{0,*}},
\end{equation}
\begin{equation}
\label{eq:part05-Cnt-micro}
\mathfrak C_{{\rm nt},*}
=
\left(\frac{\lambda_W^2\mu_W}{K_\eta^{({\rm eff})}(K_W^{({\rm eff})})^2}\right)
\left(\frac{\gamma^2\lambda_W^2\sigma}{K_UK_W^{({\rm eff})}}\right)^{E_*}
\left(\frac{\pi^2T_U}{L^2K_U}\right)^{-F_*},
\end{equation}
and
\begin{equation}
\label{eq:part05-epsilon-eta-monomial}
\epsilon_\eta=\frac{c_{\eta U}^2}{K_UK_\eta^{({\rm eff})}}.
\end{equation}

\subsection{Linear monomial-drift map}
Introduce the grouped weak-axisymmetric logarithmic drift vector
\begin{equation}
\label{eq:part05-drift-vector}
\delta\mathbf x
=
\begin{pmatrix}
\lambda_1\\ c_1\\ \gamma_1\\ \kappa_U\\ \kappa_\eta\\ \kappa_W\\ \mu_1\\ \tau_1
\end{pmatrix}
=
\begin{pmatrix}
\delta\ln\lambda_W\\
\delta\ln c_{\eta U}\\
\delta\ln\gamma\\
\delta\ln K_U\\
\delta\ln K_\eta^{({\rm eff})}\\
\delta\ln K_W^{({\rm eff})}\\
\delta\ln\mu_W\\
\delta\ln T_U
\end{pmatrix}_{\rm grp}.
\end{equation}
Then
\begin{equation}
\label{eq:part05-Mstar-drift-map}
\begin{pmatrix}
\delta\ln\mathfrak C_{{\rm tr},*}\\[2pt]
\delta\ln\mathfrak C_{{\rm nt},*}\\[2pt]
\delta\ln\epsilon_\eta
\end{pmatrix}
=
M_*\delta\mathbf x,
\end{equation}
with
\begin{equation}
\label{eq:part05-Mstar-matrix}
M_*=
\begin{pmatrix}
0 & 1+\delta_{U,*} & 1+\delta_{U,*} & -(2+\chi_{0,*}+\delta_{U,*}) & 0 & 0 & 0 & 1+\chi_{0,*}\\[4pt]
2(1+E_*) & 0 & 2E_* & F_*-E_* & -1 & -(2+E_*) & 1 & -F_*\\[4pt]
0 & 2 & 0 & -1 & -1 & 0 & 0 & 0
\end{pmatrix}.
\end{equation}
A convenient dependent minor has determinant
\begin{equation}
\label{eq:part05-Mstar-minor}
\det M_*^{(\tau,\kappa_\eta,\mu_1)}=1+\chi_{0,*}>0.
\end{equation}
Therefore \(M_*\) has rank three and a five-dimensional kernel.  The zero-defect condition is
\begin{equation}
\label{eq:part05-zero-defect-Mstar}
\Theta_1=\Xi_1=\mathcal R_1=0
\iff
M_*\delta\mathbf x=0.
\end{equation}

\claimstatus{Stage 185 is \StatusExactClosure{} in the reference coherent branch.  The monomial exponents are fixed by the branch normal form, and \(M_*\) is an exact exponent matrix.}

\section{Similarity-orbit theorem}
\label{sec:part05_weak_axisymmetric_orbit:similarity-orbit-theorem}

\subsection{Finite similarity family}
Stage 186 integrates the tangent kernel of \(M_*\) to a finite five-parameter similarity family.  Let
\begin{equation}
\label{eq:part05-free-log-scalings}
(\Lambda,C,\Gamma,U,W)
\end{equation}
be free logarithmic co-scalings of
\((\lambda_W,c_{\eta U},\gamma,K_U,K_W^{({\rm eff})})\).  Monomial preservation fixes the dependent scalings:
\begin{equation}
\label{eq:part05-Keta-orbit-law}
K_\eta^{({\rm eff})}\mapsto e^{2C-U}K_\eta^{({\rm eff})},
\end{equation}
\begin{equation}
\label{eq:part05-TU-orbit-law}
T_U\mapsto
\exp\!\left[
U-\frac{1+\delta_{U,*}}{1+\chi_{0,*}}(\Gamma+C-U)
\right]T_U,
\end{equation}
\begin{equation}
\label{eq:part05-muW-orbit-law}
\mu_W\mapsto e^{M(\Lambda,C,\Gamma,U,W)}\mu_W,
\end{equation}
where
\begin{equation}
\label{eq:part05-M-exponent-orbit-law}
\begin{aligned}
M(\Lambda,C,\Gamma,U,W)={}&
2C-U+2W-2\Lambda
-E_*(2\Gamma+2\Lambda-U-W)\\
&-F_*\frac{1+\delta_{U,*}}{1+\chi_{0,*}}(\Gamma+C-U).
\end{aligned}
\end{equation}
Call this family \(\mathcal G_*\).  It preserves \((\mathfrak C_{{\rm tr},*},\mathfrak C_{{\rm nt},*},\epsilon_\eta)\) exactly.

\subsection{Orbit--quotient closure}
Stage 187 proves the converse.  On the positive microscopic state space
\begin{equation}
\label{eq:part05-positive-state-space}
\mathcal M_+
=
\{(\lambda_W,c_{\eta U},\gamma,K_U,K_\eta^{({\rm eff})},K_W^{({\rm eff})},\mu_W,T_U)>0\},
\end{equation}
define
\begin{equation}
\label{eq:part05-invariant-map}
\mathcal I(x)
=
(\mathfrak C_{{\rm tr},*}(x),\mathfrak C_{{\rm nt},*}(x),\epsilon_\eta(x)).
\end{equation}
For two positive states \(x,\widetilde x\), equality of the three invariants is exactly
\begin{equation}
\label{eq:part05-finite-fibre-equations}
M_*\Delta\mathbf x=0,
\qquad
\Delta\mathbf x:=\ln(\widetilde x/x).
\end{equation}
Solving these finite equations gives precisely the dependent transformations \eqref{eq:part05-Keta-orbit-law}--\eqref{eq:part05-muW-orbit-law}.  Hence
\begin{equation}
\label{eq:part05-orbit-fibre-theorem}
\mathcal I(\widetilde x)=\mathcal I(x)
\iff
\widetilde x\in\mathcal G_*\cdot x.
\end{equation}
Therefore
\begin{equation}
\label{eq:part05-quotient-space}
\mathcal M_+/\mathcal G_*
\cong
(\mathbb R_{>0})^3,
\end{equation}
with quotient coordinates \eqref{eq:part05-intro-quotient-coordinates}.  The first grouped weak-axisymmetric observables are the infinitesimal motion in these quotient coordinates:
\begin{equation}
\label{eq:part05-observables-quotient}
\Theta_1=-C_{\rm tr}\delta\ln\mathfrak C_{{\rm tr},*},
\end{equation}
\begin{equation}
\label{eq:part05-Xi-quotient}
\Xi_1=A_{\rm tr}\delta\ln\mathfrak C_{{\rm tr},*}+\delta\ln\mathfrak C_{{\rm nt},*},
\end{equation}
\begin{equation}
\label{eq:part05-R-quotient}
\mathcal R_1+\Xi_1
=-\frac{\epsilon_{\eta,*}}{1-\epsilon_{\eta,*}}\delta\ln\epsilon_\eta.
\end{equation}

\subsection{Projector calculus}
Stages 191--192 then turn the quotient coordinates into an exact finite packet and an exact microscopic projector split.  With the finite log-ratio vector ordered as
\begin{equation}
\label{eq:part05-finite-log-ratio-vector}
\Delta\mathbf x
=
(\Delta_\lambda,\Delta_c,\Delta_\gamma,\Delta_U,\Delta_{K_\eta},\Delta_W,\Delta_\mu,\Delta_T)^T,
\end{equation}
the quotient packet is
\begin{equation}
\label{eq:part05-q-packet}
\mathbf q
=
(q_{\rm tr},q_{\rm nt},q_\eta)^T
=M_*\Delta\mathbf x.
\end{equation}
Choosing the dependent triple \((T_U,K_\eta^{({\rm eff})},\mu_W)\), the exact pivot block is
\begin{equation}
\label{eq:part05-pivot-block}
P_{(T,K_\eta,\mu)}
=
\begin{pmatrix}
1+\chi_{0,*} & 0 & 0\\[4pt]
-F_* & -1 & 1\\[4pt]
0 & -1 & 0
\end{pmatrix},
\qquad
\det P_{(T,K_\eta,\mu)}=1+\chi_{0,*}>0.
\end{equation}
The canonical section \(S_{(T,K_\eta,\mu)}\) is the right inverse satisfying \(M_*S_{(T,K_\eta,\mu)}=I_3\).  The exact projectors are
\begin{equation}
\label{eq:part05-projectors}
Q_{\rm quot}:=S_{(T,K_\eta,\mu)}M_*,
\qquad
O_{\rm orb}:=I_8-Q_{\rm quot}.
\end{equation}
They obey
\begin{equation}
\label{eq:part05-projector-identities}
Q_{\rm quot}^2=Q_{\rm quot},
\qquad
O_{\rm orb}^2=O_{\rm orb},
\qquad
M_*O_{\rm orb}=0,
\qquad
M_*Q_{\rm quot}=M_*.
\end{equation}
The quotient-failure piece is supported only on the dependent triple:
\begin{equation}
\label{eq:part05-failure-components}
(\Delta_T)_{\rm fail}=\frac{q_{\rm tr}}{1+\chi_{0,*}},
\qquad
(\Delta_{K_\eta})_{\rm fail}=-q_\eta,
\end{equation}
\begin{equation}
\label{eq:part05-failure-mu-component}
(\Delta_\mu)_{\rm fail}
=
\frac{F_*}{1+\chi_{0,*}}q_{\rm tr}+q_{\rm nt}-q_\eta.
\end{equation}
Thus
\begin{equation}
\label{eq:part05-orbit-lock-projector}
\mathbf q=0
\iff
M_*\Delta\mathbf x=0
\iff
Q_{\rm quot}\Delta\mathbf x=0
\iff
\Delta\mathbf x\in\ker M_*.
\end{equation}

\claimstatus{Stages 186--192 are \StatusExactClosure{} for the positive coherent reduced state space.  The word ``orbit'' here means a finite reduced monomial-preserving similarity class, not a fundamental gauge symmetry of the unreduced parent PDE.}

\section{Four-scalar packet \texorpdfstring{\((\Delta_{\rm full})\)}{(Deltafull)}}
\label{sec:part05_weak_axisymmetric_orbit:four-scalar-packet-delta-m-full}

\subsection{Packet A: conservative target and outgoing finish line}
Stages 193--197 sharpen the branch-side Packet A.  Start from the Stage 191 branch residual
\begin{equation}
\label{eq:part05-Delta-branch-full}
\Delta_{\rm branch}
=
(a_2,b_2,a_4,b_4,a_{P_0},b_{P_0},\Delta_{\rm pole},\Delta_{\rm norm}).
\end{equation}
Stage 193 freezes the conservative target surface
\begin{equation}
\label{eq:part05-conservative-target-surface}
a_2=b_2=a_4=b_4=0,
\qquad
\Delta_{\rm pole}:=\bar\nu_4-4\bar\nu_2^{\,2}=0.
\end{equation}
On the isotropic branch this is equivalent to
\begin{equation}
\label{eq:part05-one-pole-surface-D}
D_0D_4+3D_2^2=0,
\end{equation}
or, in bundle variables,
\begin{equation}
\label{eq:part05-one-pole-surface-bundle}
D_0(B_4+Z_4)=3(M+B_2+Z_2)^2.
\end{equation}
The resulting conservative carrier is
\begin{equation}
\label{eq:part05-conservative-one-pole-module}
\widehat Y_Q^{\rm cons}(\omega)
=
\frac34+\frac14\frac{1}{1-\omega^2/\Omega_Q^2}.
\end{equation}
Stage 193 also proves the scalar/geometry firewall: if weak anisotropy induces the first \(l=0\leftrightarrow l=2\) mixing, then Schur completion gives
\begin{equation}
\label{eq:part05-geometry-firewall-schur}
\mathcal D_{2,{\rm eff}}(\omega,\chi)
=
\mathcal D_2(\omega)I_3
-
\chi^2 C(\omega)\mathcal D_0(\omega)^{-1}C(\omega)^T,
\end{equation}
so
\begin{equation}
\label{eq:part05-geometry-firewall-result}
\partial_\chi\mathcal D_{2,{\rm eff}}(\omega,\chi)|_{\chi=0}=0.
\end{equation}
The scalar/geometry lane cannot linearly contaminate the grouped \(l=2\) conservative carrier.

Stage 194 attaches the outgoing \(l=2\) DtN fingerprint.  With
\begin{equation}
\label{eq:part05-z-def}
z:=\frac{a\omega}{c_s},
\end{equation}
the exact outgoing operator is
\begin{equation}
\label{eq:part05-Lambda-out}
\Lambda_2^{\rm out}(z)
=
z\frac{d}{dz}\ln h_2^{(1)}(z),
\end{equation}
with normalized response
\begin{equation}
\label{eq:part05-Y2-out-fingerprint}
\widehat Y_2^{\rm out}(z)
=
1+\frac{z^2}{9}+\frac{4z^4}{81}+i\frac{z^5}{27}+O(z^6).
\end{equation}
Matching this to the retarded grouped module fixes
\begin{equation}
\label{eq:part05-chiQ-canonical}
\chi_Q=1
\end{equation}
on the canonical compact passive/outgoing branch.  More generally, for the isotropic deformation
\begin{equation}
\label{eq:part05-DtN-deformation}
\Lambda_2^{\rm def}(z)
=
S\Lambda_2^{\rm out}(\beta z)+\Sigma_0+\Sigma_2z^2+\Sigma_4z^4+i\Sigma_5z^5+O(z^6),
\end{equation}
canonical even matching forces \(\Sigma_2\) and \(\Sigma_4\), while the outgoing scalar is
\begin{equation}
\label{eq:part05-chiQ-deformation}
\chi_Q
=
\frac{3(S\beta^5+9\Sigma_5)}{3S-\Sigma_0},
\qquad
\chi_Q-1=
\frac{3S(\beta^5-1)+\Sigma_0+27\Sigma_5}{3S-\Sigma_0}.
\end{equation}
Thus the only isotropic DtN data that can move \(\chi_Q\), while preserving the canonical even moments, are \(\beta\), \(\Sigma_0\), and \(\Sigma_5\).

Stage 195 factorizes the observable odd closure:
\begin{equation}
\label{eq:part05-source-outgoing-factorization}
\widehat m_0^{\,2}\chi_QN_Q=1.
\end{equation}
On the natural point-particle source map,
\begin{equation}
\label{eq:part05-natural-source-map}
\widehat m_0\to1,
\qquad
N_Q=\chi_Q^{-1}.
\end{equation}
Stage 196 proves that additional isotropic retarded tails beginning at \(O(\omega^7)\) do not alter any coefficient through the point-particle \(2.5\)PN order.  Stage 197 therefore gives the exact Packet-A finish line
\begin{equation}
\label{eq:part05-packetA-finish-line}
\Delta_{\rm branch}=0
\iff
\chi_Q=1
\iff
\Delta_Q:=\chi_Q-1=0.
\end{equation}

\subsection{Packet B: finite and pairwise orbit lock}
Stages 198--199 complete the orbit side.  With the same positive microscopic state vector \eqref{eq:part05-micro-state-vector}, keep the free coordinates
\begin{equation}
\label{eq:part05-free-coords}
(\lambda_W,c_{\eta U},\gamma,K_U,K_W^{({\rm eff})})
\end{equation}
and the dependent triple
\begin{equation}
\label{eq:part05-dependent-triple}
(T_U,K_\eta^{({\rm eff})},\mu_W).
\end{equation}
For a fixed invariant triple \((\mathfrak C_{{\rm tr},*},\mathfrak C_{{\rm nt},*},\epsilon_{\eta,*})\), Stage 198 solves the exact dependent orbit law and defines mismatch ratios
\begin{equation}
\label{eq:part05-mismatch-triple}
(m_T,m_K,m_\mu).
\end{equation}
The logarithmic quotient chart is
\begin{equation}
\label{eq:part05-log-chart-mismatches}
q_{\rm tr}=(1+\chi_{0,*})\ln m_T,
\qquad
q_\eta=-\ln m_K,
\qquad
q_{\rm nt}=\ln m_\mu-\ln m_K-F_*\ln m_T.
\end{equation}
Orbit lock is
\begin{equation}
\label{eq:part05-single-orbit-lock}
\Delta_{\rm orbit}=0
\iff
m_T=m_K=m_\mu=1
\iff
q_{\rm tr}=q_{\rm nt}=q_\eta=0.
\end{equation}

Stage 199 removes reference-point privilege.  For two states \(x^{(1)},x^{(2)}\), define
\begin{equation}
\label{eq:part05-pairwise-q}
\Delta_{\rm orbit}^{(2\leftarrow1)}
=
(q_{\rm tr}^{(2\leftarrow1)},q_{\rm nt}^{(2\leftarrow1)},q_\eta^{(2\leftarrow1)}).
\end{equation}
Then
\begin{equation}
\label{eq:part05-two-point-orbit-lock}
x^{(2)}\in\mathcal G_*\cdot x^{(1)}
\iff
q_{\rm tr}^{(2\leftarrow1)}=q_{\rm nt}^{(2\leftarrow1)}=q_\eta^{(2\leftarrow1)}=0.
\end{equation}
The packet is an exact additive two-point cocycle, while the multiplicative mismatch and invariant-ratio packets obey the corresponding multiplicative cocycle laws.  Thus Packet B can be attached to an orbit rather than to a selected representative.

\subsection{Reference-free full verdict packet}
Let
\begin{equation}
\label{eq:part05-target-orbit}
\mathcal O_*:=\mathcal G_*\cdot x_*
\end{equation}
be the target similarity orbit.  Stage 200 defines the additive full packet
\begin{equation}
\label{eq:part05-full-packet-additive}
\Delta_{\rm full}^{(x\mid\mathcal O_*)}
:=
\left(
\Delta_Q(x),
q_{\rm tr}^{(x\leftarrow\mathcal O_*)},
q_{\rm nt}^{(x\leftarrow\mathcal O_*)},
q_\eta^{(x\leftarrow\mathcal O_*)}
\right),
\qquad
\Delta_Q(x):=\chi_Q(x)-1.
\end{equation}
The multiplicative form is
\begin{equation}
\label{eq:part05-full-packet-multiplicative}
\mathcal V_{\rm full}^{(x\mid\mathcal O_*)}
:=
\left(
\chi_Q(x),
\mathfrak R_{\rm tr}^{(x\leftarrow\mathcal O_*)},
\mathfrak R_{\rm nt}^{(x\leftarrow\mathcal O_*)},
\mathfrak R_\eta^{(x\leftarrow\mathcal O_*)}
\right),
\end{equation}
and the mismatch form is
\begin{equation}
\label{eq:part05-full-packet-mismatch}
\mathcal M_{\rm full}^{(x\mid\mathcal O_*)}
:=
\left(
\chi_Q(x),
 m_T^{(x\leftarrow\mathcal O_*)},
 m_K^{(x\leftarrow\mathcal O_*)},
 m_\mu^{(x\leftarrow\mathcal O_*)}
\right).
\end{equation}
The exact chart conversions are
\begin{equation}
\label{eq:part05-chart-conversions-1}
\mathfrak R_{\rm tr}=e^{q_{\rm tr}}=(m_T)^{1+\chi_{0,*}},
\qquad
\mathfrak R_\eta=e^{q_\eta}=\frac{1}{m_K},
\end{equation}
\begin{equation}
\label{eq:part05-chart-conversions-2}
\mathfrak R_{\rm nt}=e^{q_{\rm nt}}=\frac{m_\mu}{m_Km_T^{F_*}}.
\end{equation}
Therefore the additive, multiplicative, and mismatch packets are exact reparameterizations of the same verdict.

The final reference-free theorem is
\begin{equation}
\label{eq:part05-reference-free-home-stretch}
\Delta_{\rm full}^{(x\mid\mathcal O_*)}=0
\iff
\Delta_{\rm branch}(x)=0
\text{ and }
x\in\mathcal O_*.
\end{equation}
Using the sharpened Packet-A and Packet-B forms, this is equivalently
\begin{equation}
\label{eq:part05-reference-free-expanded}
\Delta_{\rm full}^{(x\mid\mathcal O_*)}=0
\iff
\chi_Q(x)=1,
\qquad
q_{\rm tr}^{(x\leftarrow\mathcal O_*)}=q_{\rm nt}^{(x\leftarrow\mathcal O_*)}=q_\eta^{(x\leftarrow\mathcal O_*)}=0.
\end{equation}
This is the endpoint of Part V.  The reduced compiler algebra is finished; what remains is branch realization of the four scalar conditions.

\claimstatus{Stages 193--200 are \StatusExactClosure{} once the Packet-A isotropic grouped one-pole front end, the natural point-particle source-map reduction, and the coherent orbit quotient hierarchy are adopted.  The actual PDE may still fail any of the four final scalar conditions.  That failure would be a realization failure, not an algebraic incompleteness of the ledger.}

\section{Verification outputs and downstream use}
\label{sec:part05-verification-outputs}

The usable downstream citation packet from this chapter is as follows.
\begin{enumerate}[leftmargin=2.2em]
\item \resultanchor{MTDC-T9.1}: cite for the lower-branch drift laws, tangent-compensation theorem, weak-axisymmetric grouped signature, direct outlet map, hidden-even relation, and the loading mismatch \(\Xi_{\rm load}=N_{01}/N_0-D_{01}/D_0\).
\item \resultanchor{MTDC-T9.2}: cite for the static self-similarity theorem, wall-normalized outgoing-load theorem, square-root mixed-leg law, port co-loading theorem, effective transfer-shape collapse, and coherent support-blindness.
\item \resultanchor{MTDC-T9.3}: cite for the microscopic slippage decomposition, triangular normal form, branch-invariant monomials, rank-three drift map, similarity orbit \(\mathcal G_*\), and finite orbit--quotient closure.
\item \resultanchor{MTDC-T9.4}: cite for the branch-observable compiler, transfer-shape/outgoing-prefactor compiler, scalar no-go filters, minimal Packet A/Packet B data, and two-packet home-stretch theorem.
\item \resultanchor{MTDC-T9.5}: cite for the exact orbit/quotient projectors, isotropic grouped target surface, outgoing \(l=2\) DtN fingerprint, source-map reduction, higher-odd irrelevance, and Packet-A closure \(\Delta_{\rm branch}=0\iff\chi_Q=1\).
\item \resultanchor{MTDC-T9.6}: cite for the finite orbit law, pairwise orbit transport, orbit-representative independence, and final four-scalar verdict packet \(\Delta_{\rm full}=(\Delta_Q,q_{\rm tr},q_{\rm nt},q_\eta)\).
\end{enumerate}

\begin{verificationchecklist}
\item The off-family scalar \(\delta_\perp\) is used only as the first post-Part-IV normal coordinate; after tangent compensation, first-order isotropic bundle drift is not counted as an independent failure mode.
\item The lower compensated Family--1 drift laws are applied only under the D/N tube selection, healing-lock shell closure, and frozen \(n=5\) wall-EOS reduction.
\item The four-bundle inversion through \((\Theta_w,K_s,K_q,P_0)\) is kept distinct from a direct microscopic PDE solve; it is an algebraic compiler once those observables are known.
\item The weak-axisymmetric signature \((1,1/2,-1)\) is not confused with arbitrary grouped anisotropy.  It fixes the line \(b=3a\).
\item The no-feed-down result is first order only: grouped real \(P_2\) anisotropy can enter scalar ledgers through quadratic grouped invariants.
\item The even-preserving branch is invoked before reducing the linear grouped problem to \(\Xi_{\rm load}\).
\item Static self-similarity is stated in wall-referenced variables.  Frozen port shapes do not by themselves imply zero outgoing defect when the wall baseline moves.
\item The square-root mixed-leg law is a reduced sufficient condition under interference/hybridization rigidity, not an unconditional PDE theorem.
\item The scalar \(\Xi_1=P_1/P_0\) is the physical outgoing-prefactor slope and effective transfer-shape slope; it is not introduced as a fit parameter.
\item Coherent support-blindness applies to the direct transfer-shape defect, not to all support-dependent selected-branch quantities.
\item The triangular normal form is used with branch-adapted variables \((\Sigma_{\rm tr},\Sigma_{\rm nt},\Sigma_\eta)\); raw microscopic slopes require the stated compiler before interpretation.
\item The monomial coordinates \((\mathfrak C_{{\rm tr},*},\mathfrak C_{{\rm nt},*},\epsilon_\eta)\) are exact quotient coordinates of the reduced coherent hierarchy, not fundamental gauge invariants of the unsolved parent PDE.
\item The similarity orbit \(\mathcal G_*\) identifies reduced co-scaling fibres in the positive microscopic state space \(\mathcal M_+\); branch realization still has to be checked against the actual PDE output.
\item Packet A and Packet B are kept separate until the Stage 200 theorem: Packet A is the branch-side outgoing scalar; Packet B is the orbit-side quotient triple.
\item The equivalence \(N_Q=\chi_Q^{-1}\) is tied to the natural source-map reduction.  It should not be used when the source-map factor \(m_{\hat0}\) is still nontrivial.
\item Higher odd retarded terms are irrelevant only to the reduced point-particle \(2.5\)PN finish line; later higher-order retarded diagnostics may need them.
\item The final verdict \(\Delta_{\rm full}=0\) is an exact reduced compiler condition.  It is not an assertion that the completed moving-throat PDE has already realized the condition.
\end{verificationchecklist}

Future papers should cite Part V when replacing a long weak-axisymmetric or retarded endgame by the compact packet \((\Delta_Q,q_{\rm tr},q_{\rm nt},q_\eta)\).  They should separately state whether the actual branch has returned \(\Delta_Q=0\) and \((q_{\rm tr},q_{\rm nt},q_\eta)=(0,0,0)\), or whether those realization conditions remain open.
